\section{Conclusão}

O estudo realizado permitiu caracterizar a dinâmica do processo analisado a partir do modelo fenomenológico, das simulações em malha aberta e das representações lineares obtidas em torno do ponto de operação.
% Verificou-se que o reciclo intensifica a interação entre as unidades, fazendo com que perturbações nas entradas se propaguem pelas temperaturas e composições do sistema.
As respostas obtidas evidenciaram diferentes comportamentos dinâmicos, como ganhos positivos, ganhos negativos, sobressinal e resposta inversa.

A linearização possibilitou representar o processo por funções de transferência.
A partir dessa representação, foi possível analisar as relações específicas entre entradas e saídas, avaliar respostas ao degrau, identificar polos e zeros e interpretar características como estabilidade, resposta inversa e tempo de acomodação.
Essa etapa mostrou que o modelo linearizado é útil para a análise dinâmica local, desde que sua validade seja associada às proximidades do ponto de operação.

No domínio da frequência, a análise permitiu observar o comportamento do sistema em diferentes frequências e fundamentou o desenvolvimento de filtros capazes de atenuar ruídos sem comprometer significativamente a dinâmica do processo.
Já no tempo discreto, a escolha de $T_s = 0{,}0209 \ \mathrm{h}$ preservou as principais características dinâmicas da resposta entre $F_R$ e $x_{A1}$, e o mapeamento dos polos confirmou a estabilidade do modelo discretizado.

Dessa forma, os objetivos do trabalho foram atendidos, uma vez que o modelo fenomenológico da literatura foi simulado e analisado nos domínios do tempo contínuo, da frequência e do tempo discreto.
Os resultados reforçam que processos integrados exigem uma investigação cuidadosa, pois a dinâmica inerente a essas unidades acopladas exibe interações e comportamentos transitórios que são diretamente relevantes para a operação da planta e para o desenvolvimento de futuras estratégias de controle.
